\documentclass[10pt,a4paper]{amsart}
\usepackage[margin=19mm]{geometry}
\usepackage{amsmath,amssymb,mathtools}
\usepackage[scr=rsfs,cal=euler]{mathalfa}
\usepackage{xfrac}
\usepackage[unicode,hidelinks,bookmarks=false]{hyperref}
\newcommand{\ie}{i.e.\ }
\newcommand{\cf}{cf.\ }
\newcommand{\resp}{respectively\ }
\newcommand{\Subsection}{\subsection}
\providecommand{\nobreakdash}{\nobreak\hbox{-}}
\setlength{\emergencystretch}{2em}
\multlinegap0pt
\usepackage{amssymb}
\usepackage{xcolor}
\usepackage[shortlabels]{enumitem}
\usepackage{bm}
\usepackage{tikz}
\newtheorem{proposition}{Proposition}
\newcommand{\N}{\mathbb N}
\newcommand{\lan}{\mathcal{L}}
\newcommand{\tree}{\mathcal{T}} %nouveau mars 26
\newcommand{\val}{\mathrm{val}}
\title{A Normality Conjecture on Rational Base Number Systems}
\author{M\'elodie Andrieu, Shalom Eliahou, L\'eo Vivion}
\date{Source: arXiv:2510.11723v2 (CC BY 4.0)}
\begin{document}
\maketitle

\noindent\emph{Source and licence.} A Normality Conjecture on Rational Base Number Systems. Version arXiv:2510.11723v2 (CC BY 4.0), distributed by arXiv under the Creative Commons Attribution 4.0 International licence, http://creativecommons.org/licenses/by/4.0/, as recorded in the arXiv licence metadata for this exact version and shown on the abstract page https://arxiv.org/abs/2510.11723v2. Full licence notice: \texttt{licenses/arxiv-2510.11723-CC-BY-4.0.txt}.\par\smallskip

\noindent\textit{Editorial scope.} Counted unit: the article's proposition prop:rythm\_signature (source0.tex lines 471--472). The article's complete original proof of it is delivered with it (source0.tex lines 474--489, one \texttt{proof} environment of 16 lines). No mathematics is written, repaired or re-derived: the statement and the proof are the article's own lines, in the article's own notation, and no retained line is substituted or re-flowed.  Retained uncounted as the source-local context the proof needs: lines 306--312.  Nothing was omitted from any retained span, so the package carries no omission record.  No retained line contains a citation, so the wrapper carries no bibliography and no bibliography entry.  No retained line contains a figure environment, an image-inclusion command or drawing code, so no image is delivered and the package declares no asset dependency.  Editorial formatting only, in addition: an AMS \texttt{amsart} preamble that restates the article's own declarations and macros used by the retained lines, namely the environments \texttt{proposition} on separate counters, together with the standard packages those lines need.  The retained environments print in this document as Proposition 1 in order of appearance, whereas the article prints the same items with the numbering of its own full document.  The complete native source of the article as distributed in the arXiv e-print for this exact version is bundled unchanged as \texttt{sources/186981/source0.tex}.
\begin{proposition}\label{prop:closeness_and_extendability} Let $p > q \geq 1$ be coprime integers. Then:\\
	1. The language $\lan_{p/q}$ is prefix-closed; that is, every prefix of a word in $\lan_{p/q}$ also belongs to $\lan_{p/q}$.\\
	2. The language $\lan_{p/q}$ is right-extendable: for every $u \in \lan_{p/q}$, there exists a letter $a \in \{0, \dots, p-1\}$ such that $ua \in \lan_{p/q}$. More precisely, if $u \in \lan_{p/q}$, then 
	\[
	ua \in \lan_{p/q} \iff ua \neq \mathtt{0} \text{ and } a = -p \cdot \val_{p/q}(u) \mod{q}.
	\] 
\end{proposition}

\begin{proposition}\label{prop:rythm_signature} A breadth-first traversal \emph{of the edges} of $\tree_{p/q}$ yields the $p$-periodic sequence $(a_k)_{k\geq0}$ in $\{0,\ldots,p-1\}^{\N}$, defined by $a_0=q$ and \[a_{k+1}=a_k+q \mod p\] for every  $k \in \N$.
\end{proposition}

\begin{proof} Let $a, b \in \{0,\ldots,p-1\}$ be two consecutive edges in the breadth-first traversal of the line graph of $\tree_{p/q}$. If $a$ and $b$ originate from the same node $u$, then by Proposition~\ref{prop:closeness_and_extendability}, we have $b = a + q$.
	Otherwise, $a$ and $b$ originate from two consecutive nodes $u$ and $v$. In this case we have
	\[
	\val_{p/q}(vb) = \val_{p/q}(ua) + 1
	\quad \text{and} \quad
	\val_{p/q}(v) = \val_{p/q}(u) + 1,
	\]
	which implies
	\[
	\frac{p}{q}(\val_{p/q}(u)+1) + \frac{b}{q}
	=
	\frac{p}{q}\val_{p/q}(u) + \frac{a}{q} + 1.
	\]
	Hence $b = a + q - p$.
To conclude, note  that the first edge $a_0$ connects the node $0$, whose expansion is the empty word, to the node $1$, encoded by $q$. Therefore $a_0=q$.
\end{proof}

\end{document}
